% Gemini poster theme — A0 landscape, beamer + beamerposter.
\documentclass[final]{beamer}

\usepackage[T1]{fontenc}
\usepackage{lmodern}
\usepackage[size=custom,width=118.9,height=84.1,scale=1.0]{beamerposter}
\usetheme{gemini}
\usecolortheme{gemini}
\usepackage{graphicx}
\usepackage{booktabs}
\usepackage{tikz}

\newlength{\sepwidth}
\newlength{\colwidth}
\setlength{\sepwidth}{0.025\paperwidth}
\setlength{\colwidth}{0.3\paperwidth}
\newcommand{\separatorcolumn}{\begin{column}{\sepwidth}\end{column}}

\title{A Legible Result Communicated Over Coffee}
\author{A. Researcher \inst{1} \and B. Collaborator \inst{2} \and C. Advisor \inst{1}}
\institute[shortinst]{\inst{1} Department of Computational Something \samelineand \inst{2} University of Somewhere}

\footercontent{
	\href{https://example.org}{example.org} \hfill
	Example Conference 2026, Somewhere \hfill
	\href{mailto:a.researcher@example.org}{a.researcher@example.org}}

\begin{document}

\begin{frame}[t]
\begin{columns}[t]
\separatorcolumn

\begin{column}{\colwidth}
	\begin{block}{Introduction}
		This poster is built on the widely-forked Gemini beamerposter
		theme: a beamer frame scaled to A0 by \texttt{beamerposter}, with
		\texttt{block} environments arranged in separator-spaced columns.
		Replace this paragraph with the motivation for the work: the
		problem being solved, why it matters, and what gap the poster's
		result fills.
	\end{block}
	\begin{block}{Methods}
		\begin{itemize}
			\item Describe the approach, dataset, or apparatus here.
			\item A second bullet, for a second step of the method.
			\item A third, for whatever completes the picture.
		\end{itemize}
	\end{block}
	\begin{block}{Result Figure}
		\begin{figure}
			\centering
			\begin{tikzpicture}
				\draw[thick] (0,0) rectangle (8,6);
				\draw[thick, dashed] (0,0) -- (8,6);
				\draw[thick, dashed] (0,6) -- (8,0);
				\node at (4,3) {\Large figure placeholder};
			\end{tikzpicture}
			\caption{A placeholder figure — replace with a plot or diagram.}
		\end{figure}
	\end{block}
\end{column}

\separatorcolumn

\begin{column}{\colwidth}
	\begin{block}{Result Table}
		\centering
		\begin{tabular}{lcc}
			\toprule
			Condition & Metric A & Metric B \\
			\midrule
			Baseline  & 0.42     & 12.3     \\
			Proposed  & 0.61     & 9.7      \\
			\bottomrule
		\end{tabular}
	\end{block}
	\begin{exampleblock}{Discussion}
		Summarise what the result above means, and how it compares to
		prior work. This is also the natural place for limitations or
		threats to validity.
	\end{exampleblock}
	\begin{block}{Conclusion}
		One or two sentences stating the takeaway a passer-by should
		leave with, even if they read nothing else on the poster.
	\end{block}
\end{column}

\separatorcolumn

\begin{column}{\colwidth}
	\begin{alertblock}{Key Takeaway}
		If a reader remembers only one sentence from this poster, it
		should be this one.
	\end{alertblock}
	\begin{block}{Acknowledgements}
		Funding, collaborators, and anyone else the authors wish to
		thank.
	\end{block}
	\begin{block}{References}
		\begin{itemize}
			\item[{[1]}] A. Author, ``A Foundational Paper,''
				\emph{Journal of Examples}, 2020.
			\item[{[2]}] B. Author and C. Author, ``A Follow-up
				Study,'' \emph{Proc. of the Conference on
				Placeholders}, 2022.
		\end{itemize}
	\end{block}
\end{column}

\separatorcolumn
\end{columns}
\end{frame}

\end{document}
